\section{Build e Push da Primeira \textit{Image}}

\begin{frame}
  \frametitle{Docker Hub}
  Registro padrão para armazenar e compartilhar Docker \textit{images}.

  \vspace{0.5em}
  Duas categorias de \textit{images} confiáveis:
  \begin{itemize}
    \item \textbf{Docker Official Images (DOI)}: \textit{images} dos softwares mais populares, seguindo boas práticas e atualizadas regularmente
    \item \textbf{Docker Hardened Images (DHI)}: \textit{images}seguras e simples, desenhadas para reduzir a superfície de ataque
  \end{itemize}

  \vspace{0.5em}
  Para subir \textit{images} no Docker Hub, é necessário ter uma \textbf{conta Docker}.
\end{frame}

\begin{frame}
  \frametitle{Criando um Repositório}
  No site do Docker Hub:

  \vspace{0.5em}
  \begin{enumerate}
    \item Selecione a opção de criar repositório
    \item Preencha o nome e coloca uma descrição
    \item Defina a visibilidade — \textbf{pública} ou \textbf{privada}
  \end{enumerate}
\end{frame}

\begin{frame}
  \frametitle{Conceitos Fundamentais}
  \begin{block}{\textit{Image}}
    Pacote único com tudo o que é necessário para rodar um processo.
    Qualquer máquina com Docker pode rodar a aplicação sem instalar nada previamente.
  \end{block}

  \vspace{0.5em}
  \begin{block}{Dockerfile}
    Script com o conjunto de instruções de como buildar a image.
  \end{block}
\end{frame}

\begin{frame}
  \frametitle{Comandos de Build e Push}
  Login no Docker Hub:

  \vspace{0.3em}
  \texttt{docker login}

  \vspace{0.5em}
  Buildar a image:

  \vspace{0.3em}
  \texttt{docker build -t <DOCKER\_USERNAME>/getting-started-todo-app}

  \vspace{0.5em}
  Verificar se a image existe localmente:

  \vspace{0.3em}
  \texttt{docker image ls}

  \vspace{0.5em}
  Subir a image para o Docker Hub:

  \vspace{0.3em}
  \texttt{docker push <DOCKER\_USERNAME>/getting-started-todo-app}
\end{frame}